\section{Conclusion}
A detector that learns on the device needs a second witness, and an IoT fleet provides one. FOA lets
each device take freely whatever drift its vulnerability-diverse peers also see, and rations only what
is specific to it. It never stops learning, so it cannot be frozen; it bounds what an attacker on one
device can do and how fast; and it costs one 92-byte frame per cycle and about 22\,KB of flash on an
ESP32. On real IoT traffic with strong shared drift it was the only rule we tested that kept both
poisoning and false alarms low. What it cannot do, and what no rule can, is tell a device's own honest
change from an attack on that device; FOA turns that limit into a single, explicit rate. The next step
is the four-node ESP32 testbed, where the drift, the radio and the energy are real.

\section*{Artifacts}
Code, data scripts, experiment logs and firmware are in the project repository: the simulator and all
experiments (\texttt{experiments/foa}), the ESP32 firmware and its equivalence tests
(\texttt{firmware}), and this paper's number generator (\texttt{paper/make\_numbers.py}).

\section*{Acknowledgment}
We thank Dr.~Karnish Nasrin Ahmed Tapadar for guidance in CSD457.
